\subsection{CHARACTER FORMULA}

In this number, we write the group \(X(T)\) multiplicatively, and regard its elements as complex-valued functions on T. We assume that the element \(\rho\) of \(X(T) \otimes Q\) belongs to \(X(T)\) .

Denote by \(L^{2}(T)\) the Hilbert space of classes of square-integrable complex functions on T, and by \(\Theta(T)\) the subspace consisting of the continuous representative functions. By Spectral Theories, \(X(T)\) is an orthonormal basis of \(L^{2}(T)\) and an algebraic basis of \(\Theta(T)\) .

For \(f \in \mathrm{L}^{2}(\mathrm{T})\) and \(w \in W\) , denote by wf the element of \(\mathrm{L}^{2}(\mathrm{T})\) defined by \(wf(t) = f(w^{-1}(t))\) ; thus, for \(\lambda \in \mathrm{X}(\mathrm{T})\) we have \(w\lambda = w(\lambda)\) . Denote by \(\varepsilon : W \to \{1, -1\}\) the signature (the unique homomorphism such that \(\varepsilon(s) = -1\) for every reflection s); for \(f \in \mathrm{L}^{2}(\mathrm{T})\) , put

\[
\mathrm{J} (f) = \sum_ {w \in \mathrm{W}} \varepsilon (w) ^ {w} f.
\]

If \(\lambda \in X_{++}\) , the characters \(^w (\lambda \rho)\) are distinct; indeed, it suffices to prove that \(^w (\lambda \rho)\neq \lambda \rho\) for all \(w\neq 1\) ; but this follows from Lemma 1 (no. 1) and the fact that \(\langle \lambda \rho ,K_{\alpha}\rangle = \langle \lambda ,K_{\alpha}\rangle +1 > 0\) for every positive root \(\alpha\) . Consequently,

\[
\| \mathrm{J} (\lambda \rho) \| ^ {2} = \operatorname{Card} (\mathrm{W}) = w (\mathrm{G}).
\]

An element \(f \in \mathrm{L}^{2}(\mathrm{T})\) is said to be anti-invariant if \(^{w}f = \varepsilon(w)f\) for all \(w \in W\) (that is, if \(^{s}f = -f\) for every reflection s). We shall show that \(\frac{1}{w(\mathrm{G})}\mathrm{J}\) is the orthogonal projection of \(\mathrm{L}^{2}(\mathrm{T})\) onto the subspace of anti-invariant elements. Indeed, let \(f, f'\) be in \(\mathrm{L}^{2}(\mathrm{T})\) , with \(f'\) anti-invariant; then \(\mathrm{J}(f)\) is anti-invariant and

\[
\begin{array}{c} \langle f ^ {\prime}, \mathrm{J} (f) \rangle = \sum_ {w \in \mathrm{W}} \varepsilon (w) \langle f ^ {\prime},   ^ {w} f \rangle = \sum_ {w \in \mathrm{W}} \langle^ {w} f ^ {\prime},   ^ {w} f \rangle \\ = \sum_ {w \in \mathrm{W}} \langle f ^ {\prime},   f \rangle = w (\mathrm{G}) \langle f ^ {\prime},   f \rangle . \end{array}
\]

PROPOSITION 3. The elements \(\mathrm{J}(\lambda\rho)/\sqrt{w(\mathrm{G})}\) , for \(\lambda\in X_{++}\) , form an orthonormal basis of the subspace of anti-invariant elements of \(\mathrm{L}^{2}(\mathrm{T})\) , and an algebraic basis of the subspace of anti-invariant elements of \(\Theta(\mathrm{T})\) .

The proof is identical to that of Chap. VI, §3, no. 3, Prop. 1.

By Chap. VI, loc. cit., Prop. 2,

\[
\mathrm{J} (\rho) = \rho \prod_ {\alpha > 0} (1 - \alpha^ {- 1}) = \rho^ {- 1} \prod_ {\alpha > 0} (\alpha - 1),\tag{3}
\]

so

\[
\mathrm{J} (\rho) \overline {{\mathrm{J} (\rho)}} = \prod_ {\alpha} (\alpha - 1).\tag{4}
\]

By Cor. 2 of Th. 1 (§6, no. 2), we deduce:

Lemma 4. If \(\varphi\) and \(\psi\) are two continuous central functions on \(\mathbf{G}\) ,

\[
\int_ {\mathrm{G}} \overline {{\varphi (g)}} \psi (g) d g = \frac {1}{w (\mathrm{G})} \int_ {\mathrm{T}} \overline {{(\varphi (t) \mathrm{J} (\rho) (t))}}. (\psi (t) \mathrm{J} (\rho) (t)) d t.
\]

For all \(\lambda\in X_{++}\) , denote by \(\chi_{\lambda}\) the character of an irreducible representation of G of highest weight \(\lambda\) .

THEOREM 2 (H. Weyl). For all \(\lambda \in \mathrm{X}_{++}\) , we have \(\mathrm{J}(\rho)\chi_{\lambda}|\mathrm{T} = \mathrm{J}(\lambda \rho)\) .

The function \(\mathrm{J}(\rho)\cdot\chi_{\lambda}|T\) is anti-invariant under W, and is a linear combination with integer coefficients of elements of \(\mathrm{X(T)}\) . Thus, by Chap. VI, §3, no. 3, Prop. 1, it can be written as \(\sum_{\mu}a_{\mu}\mathrm{J}(\mu\rho)\) , where \(\mu\) belongs to \(X_{++}\) , and the \(a_{\mu}\) are integers all but finitely-many of which are zero; since \(\int_{\mathrm{G}}|\chi_{\lambda}(g)|^{2}dg=1\)

(Spectral Theories), it follows from Prop. 3 and Lemma 4 that \(\sum_{\mu}(a_{\mu})^{2} = 1\) ; thus, the \(a_{\mu}\) are all zero, except one which must be equal to 1 or \(-1\) . But the coefficient of \(\lambda\) in \(\chi_{\alpha}|T\) is equal to 1 (Th. 1), hence the coefficient of \(\lambda \rho\) in \(J(\rho)\cdot \chi_{\lambda}|T\) is equal to 1 (Chap. VI, §3, no. 3, Remark 2), which implies that \(a_{\lambda \rho} = 1\) , hence the theorem.

COROLLARY 1. With the notations of no. 3, we have in \(\mathbf{Z}[\mathrm{X(T)}]\) ,

\[
\left(\sum_ {w \in \mathrm{W}} \varepsilon (w) e ^ {w \rho}\right) \operatorname{Ch} [ \lambda ] = \sum_ {w \in \mathrm{W}} \varepsilon (w) e ^ {w \lambda} e ^ {w \rho} \text {for all} \lambda \in \mathrm{X} _ {+ +}.
\]

This follows from the theorem and commutative diagram (2) (no. 3).

COROLLARY 2. For all \(\lambda \in X_{++}\) and every regular element \(t\) of T,

\[
\chi_ {\lambda} (t) = \frac {\sum_ {w} \varepsilon (w) \lambda (w t) \rho (w t)}{\sum_ {w} \varepsilon (w) \rho (w t)}\tag{5}
\]

where the two sums are over the elements w of W.

Indeed, \(\mathrm{J}(\rho)(t)\) is non-zero, since \(t\) is regular (formula (4)).

If \(\varphi\) is a central function on \(\mathbf{G}\) , the restriction of \(\varphi\) to \(\mathrm{T}\) is invariant under \(\mathrm{W}\) , so \(\mathrm{J}(\rho). \varphi|\mathrm{T}\) is anti-invariant under \(\mathrm{W}\) . Further, by Spectral Theories and Th. 1, the family \((\chi_{\lambda})_{\lambda \in \mathrm{X}_{++}}\) is an algebraic basis of the space of central representative functions on \(\mathbf{G}\) and an orthonormal basis of the space \(\mathrm{ZL}^2(\mathrm{G})\) of classes of square-integrable central functions on \(\mathbf{G}\) .

Thus, from Prop. 3 and Th. 2 we deduce:

COROLLARY 3. The map which associates to any continuous central function \(\varphi\) on G the function \(w(\mathrm{G})^{-1/2}\mathrm{J}(\rho)(\varphi|\mathrm{T})\) induces an isomorphism from the space of central representative functions on G to the space of anti-invariant elements of \(\Theta(\mathrm{T})\) ; it extends by continuity to an isomorphism of Hilbert spaces from \(\mathrm{ZL}^2(\mathrm{G})\) to the subspace of anti-invariant elements of \(\mathrm{L}^2(\mathrm{T})\) .

COROLLARY 4. Let \(\varphi\) be a continuous central function on G. Then,

\[
\int_ {\mathrm{G}} \overline {{\chi_ {\lambda} (g)}} \varphi (g) d g = \int_ {\mathrm{T}} \overline {{\lambda (t)}} \prod_ {\alpha > 0} (1 - \alpha (t) ^ {- 1}) \varphi (t) d t = \int_ {\mathrm{T}} \overline {{\lambda \rho (t)}}. \varphi (t) \mathrm{J} (\rho) (t) d t.
\]

Indeed, by Lemma 4 and Th. 2,

\[
\begin{array}{c} \int_ {\mathrm{G}} \overline {{\chi_ {\lambda} (g)}} \varphi (g) d g = \frac {1}{w (\mathrm{G})} \int_ {\mathrm{T}} \overline {{\chi_ {\lambda} (t) \mathrm{J} (\rho) (t)}} (\varphi (t) \mathrm{J} (\rho) (t)) d t \\ = \frac {1}{w (\mathrm{G})} \int_ {\mathrm{T}} \overline {{\mathrm{J} (\lambda \rho) (t)}} \varphi (t) \mathrm{J} (\rho) (t) d t. \end{array}
\]

But the function \(t \mapsto \varphi(t)J(\rho)(t)\) is anti-invariant and \(\frac{1}{w(G)}J(\lambda\rho)\) is the orthogonal projection of \(\lambda\rho\) onto the subspace of anti-invariant elements of \(L^{2}(T)\) , so

\[
\frac {1}{w (\mathrm{G})} \int_ {\mathrm{T}} \overline {{\mathrm{J} (\lambda \rho) (t)}} \varphi (t) \mathrm{J} (\rho) (t) d t = \int_ {\mathrm{T}} \overline {{\lambda \rho (t)}} \varphi (t) \mathrm{J} (\rho) (t) d t;
\]

finally, by formula (3), we have \(\overline{\rho(t)}\mathrm{J}(\rho)(t) = \prod_{\alpha >0}(1 - \alpha (t)^{-1})\) , hence the corollary.

Remarks. 1) For all \(w \in \mathrm{W}\) , put \(\rho_w = {}^w\rho /\rho\) ; then

\[
\sum_ {w} \varepsilon (w) \rho_ {w} = \prod_ {\alpha > 0} (1 - \alpha^ {- 1}) = \rho^ {- 2} \prod_ {\alpha > 0} (\alpha - 1).\tag{6}
\]

If \(t\) is a regular element of \(\mathrm{T}\) , we deduce from (5) that

\[
\chi_ {\lambda} (t) = \frac {\sum_ {w} \varepsilon (w) ^ {w} \lambda (t) \rho_ {w} (t)}{\sum_ {w} \varepsilon (w) \rho_ {w} (t)} = \frac {\sum_ {w} \varepsilon (w) ^ {w} \lambda (t) \rho_ {w} (t)}{\prod_ {\alpha > 0} (1 - \alpha (t) ^ {- 1})}.\tag{7}
\]

Note that \(\rho_{w}\) is a linear combination of roots with integer coefficients, and hence belongs to X(T) even if we do not assume that \(\rho \in \mathrm{X}(T)\) . It follows that formula (7) is valid without the assumption that \(\rho \in \mathrm{X}(T)\) : indeed, to prove this we replace G by a suitable connected covering, and are then reduced to Cor. 2.

\begin{enumerate}
\def\labelenumi{\arabic{enumi})}
\setcounter{enumi}{1}
\item
  Similarly, the first equality of Cor. 4 remains valid without the assumption that \(\rho \in \mathrm{X}(\mathrm{T})\) .
\item
  We can deduce Th. 2 from the analogous infinitesimal statement (Chap. VIII, §9, no. 1, Th. 1); the same is true for Th. 3 of the next number (which is the analogue of Th. 2 of loc. cit., no. 2).
\end{enumerate}

